% =====================================================================
%  Supporting Information for
%  "Evaluating generative antibody-drug conjugate linkers"
%  Compile:  latexmk -pdf -bibtex supplementary.tex
%  (standalone; reuses references.bib)
% =====================================================================
\documentclass[10pt]{article}

\usepackage[letterpaper,margin=2.2cm]{geometry}
\usepackage[T1]{fontenc}
\usepackage[utf8]{inputenc}
\usepackage{lmodern}
\usepackage{microtype}
\usepackage{amsmath,amssymb}
\usepackage[version=4]{mhchem}
\usepackage{graphicx}
\graphicspath{{../figures/}}
\usepackage{xcolor}
\usepackage{colortbl}
\usepackage{booktabs}
\usepackage{array}
\usepackage{enumitem}
\usepackage{caption}
\usepackage[numbers,sort&compress]{natbib}
\usepackage[hidelinks]{hyperref}

\definecolor{ink}{HTML}{17314A}
\definecolor{accent}{HTML}{0072B2}
\definecolor{tint}{HTML}{EAF3FA}
\hypersetup{colorlinks=true,linkcolor=accent,citecolor=accent,urlcolor=accent}
\captionsetup{font=small,labelfont={bf,color=accent},labelsep=period}

% Number SI floats/sections as S1, S2, ...
\renewcommand{\thesection}{S\arabic{section}}
\renewcommand{\thetable}{S\arabic{table}}
\renewcommand{\thefigure}{S\arabic{figure}}
\renewcommand{\theequation}{S\arabic{equation}}

\title{\textbf{\color{ink}Supporting Information}\\[0.3em]
\large Evaluating generative antibody--drug conjugate linkers:\\
chemical validity, covalent dynamics, and an FDA-anchored control}
\author{Ga\"el Coulombe \and Hazem Mslati}
\date{}

\begin{document}
\maketitle

\noindent
This Supporting Information documents the generation pipeline (dataset construction, fine-tuning recipe,
sampling sweeps) and the chemical-validity gate. The generation work is \emph{prior} to the paper's
contribution (the evaluation result); it is reported here for reproducibility. All numbers are reported
from disk-verified pipeline outputs. Section references (\S) point to the main text.

\tableofcontents
\vspace{1em}

% =====================================================================
\section{Dataset construction}
\label{si:data}

\paragraph{Payload-anchored decomposition.}
Each ADC in the source dataset\cite{mslati2026} provides the full conjugate SMILES, the isolated payload
SMILES, and a free-text conjugation-site annotation. Earlier pipelines guessed the linker/payload
boundary with reaction SMARTS, which covered only $\approx$72\,\% of records and failed silently on
non-canonical chemistries. We instead \emph{anchor} the boundary: the provided payload is located as a
subgraph of the full conjugate, and the single cross-boundary bond defines the linker/payload cut. The
dataset stores the payload in its post-conjugation form, so 49/53 cysteine ADCs contain the payload as an
exact substructure; an MCS fallback (RDKit \texttt{rdFMCS}, \texttt{bondCompare=CompareAny}, $\geq$70\,\%
coverage) recovers the remainder. Coverage rises to \textbf{98.1\,\%} (52/53). When several substitution
sets match, the one with the fewest junction atoms is chosen, and exactly one cross-boundary bond is
required.

\paragraph{Deduplication and leakage control.}
Deduplication at the \emph{canonical linker--payload (L/P)} level --- rather than by source record UUID
--- caught four real train/validation leaks that a UUID-only check had masked. The cysteine-only
fine-tuning set is \textbf{43 unique L/Ps} (36 in train), with \textbf{0} cross-split L/P leakage. The
multi-site set extends this to \textbf{73 unique L/Ps} across three conjugation chemistries (cysteine,
lysine, click).

\paragraph{Conformer augmentation.}
Each L/P is embedded in 3D with ETKDGv3 and MMFF94 optimization (UFF fallback), generating up to 20
diverse conformers per L/P (RMS pruning, minimum-diversity 0.5\,\AA). DiffLinker consumes a fragment file
(payload + maleimide stub) and a link file (the spacer to be regenerated), with anchor indices marking the
attachment atoms. The augmented cysteine-only set is 720 train / 140 validation examples; the multi-site
set is larger (Table~\ref{tab:datasets}). Augmentation increases \emph{geometric} diversity only; it does
not add chemical (topological) diversity --- 20 conformers of one molecule remain one molecule, which is
the origin of the ``data wall'' discussed below.

\begin{table}[t]
\caption{DiffLinker training sets. ``L/Ps'' = unique linker--payloads; ``examples'' = conformer-augmented
fragment/link pairs after the train/validation split.}
\label{tab:datasets}
\centering\small
\begin{tabular}{@{}lrrl@{}}
\toprule
\rowcolor{tint}Set & unique L/Ps & examples (train/val) & chemistries \\
\midrule
Cysteine-only (dedup) & 43 & 720 / 140 & cysteine \\
Multi-site            & 73 & 1240 / 202 & cysteine, lysine, click \\
\bottomrule
\end{tabular}
\end{table}

% =====================================================================
\section{Fine-tuning recipe and numerical stability}
\label{si:finetune}

DiffLinker\cite{igashov2024} ($E(n)$-equivariant diffusion, pretrained on GEOM-Drugs\cite{axelrod2022})
was fine-tuned on a single RTX~4080 (16\,GB). Hyperparameters: learning rate $1\times10^{-5}$, batch
size 4, gradient clipping 1.0, \emph{no} learning-rate warmup, \emph{no} optimizer reset (the GEOM Adam
moment estimates are restored), early stopping (patience 20) on the denoising validation loss. Best
validation loss was 0.1159 for the cysteine-only fine-tune and 0.1143 for the multi-site fine-tune.

\paragraph{The isfinite guard.}
Large ADCs ($\gtrsim$95 heavy atoms) at high diffusion timesteps occasionally produce a non-finite
forward loss in single precision (a physics/numerics overflow, $\approx$2\,\% of batches). A pre-backward
\texttt{isfinite} guard that skips the offending batch --- combined with the restored Adam moments ---
prevents a NaN cascade. In a controlled ablation (5 runs $\times$ 200 steps), the guard gave 0/5 NaN
cascades and 5/5 convergence; without it, 4/5 runs diverged. Resetting the optimizer or adding warmup
\emph{worsened} stability, because the restored GEOM moments actively pull weights out of the overflow
region.

\paragraph{Denoising loss is not generation quality.}
The validation loss is a diffusion \emph{denoising} objective and is \emph{not} a direct measure of
generated-molecule quality. We therefore report generation quality with the operational usable-3D yield
and Wilson 95\,\% confidence intervals (Section~\ref{si:sweeps}), not with validation loss.

% =====================================================================
\section{The two tradeoffs and the Pareto frontier}
\label{si:pareto}

Two intrinsic tradeoffs constrain the generator; neither is removed by any easy lever, so the operating
point can only be \emph{moved along} a Pareto frontier.

\paragraph{Chemistry vs.\ connectivity (the size axis).}
Within one model, forcing larger linkers increases ADC chemistry but decreases geometric cohesion. As the
forced linker size rises from 15 to 60, single-fragment connectivity falls ($\approx$30\,\%
$\rightarrow$ $\approx$4\,\%) while the Val-Cit and urea motif rates rise (to 64\,\% and 81\,\%
respectively; Table~1, main text). The operational size is chosen near the knee (20--25).

\paragraph{Specificity vs.\ diversity (the dataset axis).}
Broadening the training distribution (multi-site) improves connectivity but dilutes the dominant
cysteine--Val-Cit signal toward the population mean. The two effects trade off monotonically with the
\emph{cysteine sampling fraction}, which therefore acts as an operational slider along the frontier
(Table~\ref{tab:pareto}). No setting maximizes both connectivity and target chemistry.

\begin{table}[t]
\caption{The Pareto frontier: cysteine sampling fraction vs.\ connectivity and target chemistry, at the
operational size (20). Connectivity and Val-Cit rates move in opposition; no configuration is dominant.}
\label{tab:pareto}
\centering\small
\begin{tabular}{@{}lrrr@{}}
\toprule
\rowcolor{tint}Configuration & \% cysteine & native-conn.\ (\%) & Val-Cit (\%) \\
\midrule
Multi-site (natural)   & 66  & 45 & 10 \\
Curriculum ($3\times$) & 85  & 23 & 33 \\
Cysteine-only          & 100 & 28 & 23 \\
\bottomrule
\end{tabular}
\end{table}

% =====================================================================
\section{Sampling sweeps and the $n{=}500$ re-baseline}
\label{si:sweeps}

\paragraph{Usable-3D definition.}
A generated molecule is \emph{usable-3D} if it is (a)~natively connected (a single fragment under
distance-based bond perception, with \emph{no} post-hoc minimum-spanning-tree closure), (b)~MMFF94-valid
(radical-neutralized, sanitizable, and minimizable to a finite energy while remaining a single fragment),
and (c)~ADC-motif-bearing (Val-Cit or urea). This is the gate of \S\,4.2; it is the strict 3D criterion
required for the downstream covalent build, which a 2D-only screen would not satisfy.

\paragraph{Re-baseline.}
Early single-checkpoint comparisons at $n{=}100$ suggested a $\approx$1.5--2$\times$ yield gain from 50
extra fine-tuning epochs. Re-running at $n{=}500$ with Wilson 95\,\% confidence intervals
(Table~\ref{tab:rebaseline}) erased it: a clean $\lambda{=}0$ control --- in which the additional
connectivity penalty is a mechanical no-op --- regressed to baseline rather than matching the $\lambda{=}0.01$
arm, with \emph{disjoint} confidence intervals. The only difference between the two runs is
GPU-nondeterministic training-trajectory variance. \textbf{The apparent gain was a single
favourable-variance draw.} The operational yield is $\approx$9--10 usable-3D candidates per 100. The
methodological consequence: single-checkpoint comparisons at $n{<}500$ cannot separate variance from
signal; multi-seed evaluation with confidence intervals at $n{\geq}500$ is required.

\begin{table}[t]
\caption{$n{=}500$ re-baseline (3D-strict), Wilson 95\,\% CI. The $\lambda{=}0.01$ arm and the
$\lambda{=}0$ control have disjoint CIs despite the penalty being a no-op in the control --- the
difference is run-to-run training variance, not a real effect. Curriculum, multi-site, and the
$\lambda{=}0$ control are mutually overlapping.}
\label{tab:rebaseline}
\centering\small
\setlength{\tabcolsep}{5pt}
\begin{tabular}{@{}lrrrr@{}}
\toprule
\rowcolor{tint}Config ($n{=}500$) & conn.\ \% & Val-Cit \% & usable-3D \% & 95\,\% CI \\
\midrule
Curriculum (baseline)          & 20.8 & 6.6 & 8.6  & [6.4,\,11.4] \\
Multi-site (base)              & 40.0 & 4.4 & 4.8  & [3.2,\,7.0] \\
$\lambda{=}0.01$ (prev.\ best) & 33.2 & 9.4 & 10.2 & [7.8,\,13.2] \\
$\lambda{=}0$ control          & 15.4 & 4.4 & 4.6  & [3.1,\,6.8] \\
\bottomrule
\end{tabular}
\end{table}

\paragraph{Novelty.}
Across 1000 generations there were \textbf{zero} near-copies (Tanimoto $>0.85$ to any training L/P), and
the median Tanimoto to the training set was 0.44--0.59 --- novel molecules on a shared payload scaffold,
not memorization.

% =====================================================================
\section{Chemical-validity gate specification}
\label{si:gate}

The gate runs fail-closed in two levels (\S\,3.2).

\paragraph{Level A --- ligand (RDKit\cite{rdkit}).}
\begin{enumerate}[leftmargin=1.4em,itemsep=1pt,topsep=2pt]
\item Sanitizable (valence-consistent Kekul\'e/aromaticity perception).
\item Single connected component (rejects the diffusion fragmentation that dominates raw output).
\item No undefined stereocenters on linker or payload (the coordinate$\rightarrow$SDF conversion strips
stereochemistry; canonical payload stereochemistry is reimposed by CIP-driven iterative chiral-tag
assignment and verified from the 3D coordinates).
\item Correct maleimide vs.\ thiosuccinimide oxidation state (open maleimide \ce{O=C1C=CC(=O)N1}
distinguished from the conjugated thiosuccinimide).
\item Basic drug-likeness (molecular weight, logP, rotatable bonds, TPSA within ranges).
\end{enumerate}

\paragraph{Level B --- conjugate topology.}
A text-level parse of the assembled GROMACS topology enforces a per-atom maximum-valence table (carbon
$\leq$4, sulfur $\leq$2 for thioether, etc.). This is the guard that auto-catches the over-valent
conjugation carbon of \S\,4.4 and the hypervalent-sulfur defect (a retained cysteine thiol proton) ---
both of which pass a distance-only check, because a harmonic bond holds its length regardless of valence.
The lesson (Table~2, main text): validate the reacting atom's bond \emph{count and type}, not just the
new bond's \emph{length}.

% =====================================================================
\section{Why the AMM potency oracle does not rank linkers}
\label{si:amm}

The AMM classifier\cite{mslati2026} predicts whether an ADC is cytotoxic below a 10\,nM threshold from a
2D conjugate representation plus biological context. It does not rank the generated linkers, for two
compounding reasons. (i)~\textbf{Label saturation}: a potent payload (MMAE) on an amplified antigen (HER2)
at a 10\,nM threshold is predicted active for essentially any reasonable linker. (ii)~\textbf{No
applicability-domain check}: after a silent checkpoint-loading defect was corrected (a vestigial untrained
head had to be dropped and only the attention keys remapped; diagnosed via BatchNorm tracking counters),
the genuine model still scores chemically irrelevant ``junk'' control molecules (aspirin, caffeine,
ethanol) as active --- their molecular features are out-of-distribution, propagating to extreme logits.
AMM is therefore usable only as a coarse plausibility filter (real-vs-junk), never as a fine-grained
ranker. The discriminating signal --- linker-driven plasma stability --- was never in its training data.

% =====================================================================
{\small
\bibliographystyle{unsrtnat}
\bibliography{references}
}

\end{document}
